\section{Bell--Gould prime-power filters and mixed integrals}
\label{sec:touchard-mixed}

The purpose of this section is twofold.  First we isolate an exact
integer-valued-polynomial mechanism which annihilates the Bell and Gould
functionals simultaneously modulo an arbitrary integer.  The boundary term
in the Gould functional is retained throughout; omitting it makes the
prime-power lift invalid.  Second, we match the resulting reciprocal
integrals with logarithmic moments.  The matching places the answer in the
pure \(E\)-value direction \(e^\lambda\Ein(\lambda)\), where
Theorem~\ref{thm:pure-exp-base} turns the construction into an unconditional
algebraic-independence theorem.

\subsection{The Bell--Gould adjoint calculus}

Let \((\mathsf B_n)_{n\geq0}\) be the Bell sequence and let
\((\mathsf G_n)_{n\geq0}\) be the Gould sequence, normalized by
\[
 \mathsf B_0=1,
 \qquad
 \mathsf B_{n+1}=\sum_{k=0}^n\binom nk\mathsf B_k,
\]
and
\[
 \mathsf G_0=0,
 \quad \mathsf G_1=1,
 \qquad
 \mathsf G_{n+1}=\sum_{k=0}^n\binom nk\mathsf G_k
 \quad(n\geq1).
\]
Define the two integral linear functionals on \(\Z[x]\) by
\[
 \mathfrak B(x^n)=\mathsf B_n,
 \qquad
 \mathfrak G(x^n)=\mathsf G_n.
\]
For \(r\geq1\), write
\((x)_r=x(x-1)\cdots(x-r+1)\), and put
\begin{equation}\label{eq:touchard-theta}
 \Theta_r(f)=
 \sum_{j=0}^{r-1}(-1)^{r-1-j}(r-1-j)!f(j).
\end{equation}

\begin{lemma}[Bell--Gould adjoints]
\label{lem:bell-gould-adjoints}
For every \(r\geq1\) and \(f\in\Z[x]\),
\begin{align}
 \mathfrak B\bigl((x)_rf(x)\bigr)
   &=\mathfrak B\bigl(f(x+r)\bigr),
   \label{eq:bell-adjoint}\\
 \mathfrak G\bigl((x)_rf(x)\bigr)
   &=\mathfrak G\bigl(f(x+r)\bigr)+\Theta_r(f).
   \label{eq:gould-adjoint}
\end{align}
\end{lemma}

\begin{proof}
For \(r=1\), the two identities are precisely the defining binomial
recurrences, with the exceptional initial value \(\mathsf G_1=1\) producing
the term \(f(0)\) in the second identity.  Iteration proves
\eqref{eq:bell-adjoint}.  For the Gould identity, suppose it holds for
\(r\).  Apply the case \(r=1\) to
\(h(x)=(x-1)_rf(x)\), and then apply the induction hypothesis to
\(f(x+1)\).  The new boundary contribution is
\[
 (-1)^rr!f(0)+
 \sum_{j=0}^{r-1}(-1)^{r-1-j}(r-1-j)!f(j+1),
\]
which is \(\Theta_{r+1}(f)\).  This completes the induction.
\end{proof}

Let
\begin{equation}\label{eq:signed-touchard}
 \tau_n(x)=\sum_{k=0}^n
 \left\{\begin{matrix}n\\k\end{matrix}\right\}(-x)^k
\end{equation}
be the signed Touchard polynomial.  The elementary reciprocal integral is
\begin{equation}\label{eq:touchard-remainder}
 \mathscr R_n
 :=\int_0^\infty\frac{e^{-x}\tau_n(x)}{1+x}\,dx
 =\mathsf B_n\delta-\mathsf G_n.
\end{equation}
Indeed, the exponential generating function of the left side is
\[
 \int_0^\infty\frac{e^{-e^z x}}{1+x}\,dx
 =e^{e^z}E_1(e^z).
\]
The function \(\delta\exp(e^z-1)\) minus this generating function has
derivative equal to itself multiplied by \(e^z\), plus \(1\), and hence has coefficients
\((\mathsf G_n)\) with the displayed normalization.

\subsection{A boundary-aware prime-power lift}

Fix a prime \(p\).  Put
\[
 D_p(x)=x^p-x.
\]
There is a \(G_p\in\Z[x]\) such that
\begin{equation}\label{eq:dp-falling}
 D_p(x)=(x)_p+pG_p(x).
\end{equation}
To see this, subtract the two monic degree-\(p\) polynomials.  The difference
has degree at most \(p-1\) and vanishes at all \(p\) elements of
\(\mathbf F_p\), so every coefficient is divisible by \(p\).

For \(r\geq1\), set
\[
 m_{p,r}=\min\{m\geq0:v_p(m!)\geq r\}.
\]
Choose \(a\in\Z\), put \(W_{p,r,a}(x)=(x-a)_{m_{p,r}}\), and define
\begin{equation}\label{eq:touchard-lift-operator}
 \mathcal L_p[P](x)=D_p(x)P(x)-P(x+p).
\end{equation}

For a polynomial \(P\in\Z[x]\), use the abbreviations
\begin{align*}
 \mathrm{FD}_r(P)&:\Longleftrightarrow
      P(k)\in p^r\Z\quad(k\in\Z),\\
 \mathrm{AN}_s(P)&:\Longleftrightarrow
      \mathfrak B(Pf),\mathfrak G(Pf)\in p^s\Z
      \quad(f\in\Z[x]).
\end{align*}

\begin{lemma}[One-step lift]
\label{lem:touchard-one-step}
Suppose that \(1\leq s<r\) and that \(P\) satisfies
\(\mathrm{FD}_r(P)\) and \(\mathrm{AN}_s(P)\).  Then
\(Q=\mathcal L_p[P]\) satisfies
\(\mathrm{FD}_r(Q)\) and \(\mathrm{AN}_{s+1}(Q)\).
\end{lemma}

\begin{proof}
First note the shifted-annihilator consequence of
Lemma~\ref{lem:bell-gould-adjoints}.  If \(r\geq s\), then for every
\(g\in\Z[x]\),
\begin{align*}
 \mathfrak B(P(x+p)g(x))
 &=\mathfrak B(P(x)(x)_pg(x-p)),\\
 \mathfrak G(P(x+p)g(x))
 &=\mathfrak G(P(x)(x)_pg(x-p))
   -\Theta_p(P(x)g(x-p)).
\end{align*}
The first terms on the right are divisible by \(p^s\), and every summand in
the boundary term is divisible by \(p^r\).  Thus both shifted functionals
are divisible by \(p^s\).

The fixed-divisor assertion for \(Q\) follows immediately from
\eqref{eq:touchard-lift-operator}.  Write
\[
 f(x+p)-f(x)=p\Delta_pf(x),
 \qquad \Delta_pf\in\Z[x].
\]
Using \eqref{eq:dp-falling} and
Lemma~\ref{lem:bell-gould-adjoints} gives
\begin{align*}
 \mathfrak B(Qf)
 &=p\mathfrak B(P(x+p)\Delta_pf)
   +p\mathfrak B(PG_pf),\\
 \mathfrak G(Qf)
 &=p\mathfrak G(P(x+p)\Delta_pf)
   +p\mathfrak G(PG_pf)+\Theta_p(Pf).
\end{align*}
The first two terms in each line are divisible by \(p^{s+1}\).  The last
term in the second line is divisible by \(p^r\), hence by \(p^{s+1}\).
\end{proof}

\begin{theorem}[Common Bell--Gould prime-power annihilator]
\label{thm:bell-gould-prime-power}
For every prime \(p\), \(r\geq1\), and \(a\in\Z\), the polynomial
\begin{equation}\label{eq:prime-power-annihilator}
 S_{p,r}^{(a)}(x)=
 \mathcal L_p^{\,r-1}
 \left[W_{p,r,a}(x)(D_p(x)-1)\right]
\end{equation}
is monic and primitive, has degree \(m_{p,r}+rp\), and satisfies
\begin{equation}\label{eq:prime-power-universal}
 p^r\mid\mathfrak B(S_{p,r}^{(a)}f),
 \qquad
 p^r\mid\mathfrak G(S_{p,r}^{(a)}f)
 \quad(f\in\Z[x]).
\end{equation}
\end{theorem}

\begin{proof}
Put \(P_0=W_{p,r,a}(D_p-1)\).  Since
\[
 W_{p,r,a}(k)=m_{p,r}!\binom{k-a}{m_{p,r}}
 \quad(k\in\Z),
\]
the polynomial \(P_0\) satisfies \(\mathrm{FD}_r(P_0)\).  For
\(h=W_{p,r,a}f\), equations
\eqref{eq:bell-adjoint}--\eqref{eq:dp-falling} give
\begin{align*}
 \mathfrak B(P_0f)
 &=\mathfrak B(h(x+p)-h(x))+p\mathfrak B(G_ph),\\
 \mathfrak G(P_0f)
 &=\mathfrak G(h(x+p)-h(x))+p\mathfrak G(G_ph)
   +\Theta_p(h).
\end{align*}
The difference \(h(x+p)-h(x)\) belongs to \(p\Z[x]\), while every value
\(h(j)\) occurring in \(\Theta_p(h)\) is divisible by \(p^r\).  Hence
\(P_0\) satisfies \(\mathrm{AN}_1(P_0)\).  Iterating
Lemma~\ref{lem:touchard-one-step} \(r-1\) times proves
\eqref{eq:prime-power-universal}.

The starting polynomial is monic of degree \(m_{p,r}+p\).
The operator \(\mathcal L_p\) preserves monicity and raises the degree by
\(p\), since its first summand has degree \(p+\deg P\) and its second has
degree \(\deg P\).  This proves the degree formula.  A monic integral
polynomial is primitive.
\end{proof}

The carrier \(W_{p,r,a}\) has the smallest possible degree among primitive
integer polynomials whose fixed divisor is divisible by \(p^r\): the fixed
divisor of a primitive degree-\(d\) polynomial divides \(d!\), while
\(W_{p,r,a}\) attains \(d=m_{p,r}\); see
\cite{PrasadRajkumarReddy2019}.  This is an optimality statement for the
fixed-divisor carrier, not for all possible simultaneous Bell--Gould
annihilators.

\subsection{Coefficientwise CRT for an arbitrary modulus}

\begin{theorem}[Arbitrary-modulus Bell--Gould filter]
\label{thm:bell-gould-arbitrary-modulus}
Let
\[
 M=\prod_{p\mid M}p^{r_p}\geq2,
 \qquad
 d_p=m_{p,r_p}+r_pp,
 \qquad
 D_M=\max_{p\mid M}d_p.
\]
There is a monic primitive \(C_M\in\Z[x]\), of degree \(D_M\), such that
\begin{equation}\label{eq:arbitrary-modulus-universal}
 M\mid\mathfrak B(C_Mf),
 \qquad
 M\mid\mathfrak G(C_Mf)
 \quad(f\in\Z[x]).
\end{equation}
Every nonleading coefficient of \(C_M\) may be chosen to have absolute
value at most \(M/2\).
\end{theorem}

\begin{proof}
For every \(p\mid M\), choose an integer \(a_p\) and put
\[
 T_p(x)=x^{D_M-d_p}S_{p,r_p}^{(a_p)}(x).
\]
All \(T_p\) are monic of degree \(D_M\).  Apply the Chinese remainder
theorem independently to their nonleading coefficients and set the leading
coefficient equal to \(1\).  This gives
\[
 C_M\equiv T_p\pmod{p^{r_p}}\qquad(p\mid M),
\]
with centered representatives bounded by \(M/2\).  Multiplication by a
power of \(x\) preserves \eqref{eq:prime-power-universal}; moreover
\(C_M-T_p\in p^{r_p}\Z[x]\).  Since the two functionals take integral
values on \(\Z[x]\), both values in
\eqref{eq:arbitrary-modulus-universal} are divisible by every
\(p^{r_p}\), and hence by \(M\).
\end{proof}

Let (E) denote the forward shift in (n),
\[
 (Ea)_n=a_{n+1},
\]
acting coefficientwise on polynomial-valued sequences.

\begin{corollary}[Factorial specialization]
\label{cor:bell-gould-factorial}
For \(M=N!\), \(N\geq2\), one may take a monic common filter of degree
\(D_N\) satisfying
\begin{equation}\label{eq:factorial-degree-bound}
 3N-2\lfloor\log_2N\rfloor-3
 \leq D_N\leq3N.
\end{equation}
Consequently, for every \(n\geq0\), there are integers \(P_{N,n},Q_{N,n}\)
such that
\begin{equation}\label{eq:factorial-touchard-integral}
 \frac1{N!}\int_0^\infty
 \frac{e^{-x}[C_{N!}(E)\tau_n](x)}{1+x}\,dx
 =Q_{N,n}\delta-P_{N,n}.
\end{equation}
\end{corollary}

\begin{proof}
If \(r_p=v_p(N!)\), then
\[
 m_{p,r_p}=p\lfloor N/p\rfloor,
 \qquad
 d_p=p\lfloor N/p\rfloor+pv_p(N!).
\]
The quotient of \(N!\) by
\((p\lfloor N/p\rfloor)!\) is a \(p\)-adic unit, and deleting the last
multiple of \(p\) lowers the valuation, proving the first equality.  Also
\[
 d_p\leq N+p\sum_{j\geq1}\frac{N}{p^j}\leq3N.
\]
For \(p=2\), Legendre's formula
\(v_2(N!)=N-s_2(N)\), together with
\(s_2(N)\leq\lfloor\log_2N\rfloor+1\), gives the lower bound in
\eqref{eq:factorial-degree-bound}.  Finally apply
\eqref{eq:touchard-remainder} and
Theorem~\ref{thm:bell-gould-arbitrary-modulus} with \(f=x^n\).
\end{proof}

The congruence tradition begins with Touchard's classical congruence for
Bell numbers and its polynomial variants
\cite{Touchard1933,SunZagier2011}; the adjoint language belongs to the
classical umbral calculus \cite{Roman1984}, and the companion Gould
recurrence was studied systematically by Gould and Quaintance
\cite{GouldQuaintance2007}.  The preceding proof records a particular
boundary-aware simultaneous lift and its arbitrary-modulus CRT packaging.
No claim of historical priority for that exact packaging is made without a
separate search through the older congruence and umbral literature.

\subsection{Matched mixed moments}

Retain the pure exponential field
\[
 \Eexp=\Qbar(e^\alpha:\alpha\in\Qbar).
\]
For \(U(t)=\sum_ku_kt^k\in\Qbar[t]\), define
\begin{equation}\label{eq:matched-beta-r}
 \beta(U)=\sum_ku_kk!,
 \qquad
 R(U)=\sum_ku_kk!H_k.
\end{equation}
For \(\lambda\in\Qbar_{>0}\) and
\(V(t)=\sum_kv_kt^k\in\Qbar[t]\), put
\begin{equation}\label{eq:matched-s}
 S_\lambda(V)=
 \sum_{k\geq1}v_k\sum_{j=0}^{k-1}
 (-\lambda)^j(k-1-j)!.
\end{equation}

\begin{theorem}[Matched mixed-integral independence]
\label{thm:matched-mixed-integrals}
Let \(\lambda_1,\ldots,\lambda_s\) be distinct positive algebraic numbers.
For each \(i\), choose \(U_i,V_i\in\Qbar[t]\) such that
\begin{equation}\label{eq:matched-condition}
 V_i(-\lambda_i)=-\beta(U_i)\ne0,
\end{equation}
and define
\begin{equation}\label{eq:matched-integral}
 \mathscr I_i=
 e^{\lambda_i}\int_0^\infty e^{-t}U_i(t)
       \log\!\left(\frac{t}{\lambda_i}\right)dt
 +\int_0^\infty\frac{e^{-t}V_i(t)}{t+\lambda_i}\,dt.
\end{equation}
Then \(\mathscr I_1,\ldots,\mathscr I_s\) are algebraically independent
over \(\Eexp\).
\end{theorem}

\begin{proof}
The logarithmic moments are
\begin{equation}\label{eq:matched-log-moment}
 \int_0^\infty e^{-t}t^k\log(t/\lambda)\,dt
 =k!(H_k-\gamma-\log\lambda).
\end{equation}
If
\(J_k(\lambda)=\int_0^\infty e^{-t}t^k/(t+\lambda)\,dt\), then
\[
 J_0(\lambda)=e^\lambda E_1(\lambda),
 \qquad
 J_k(\lambda)=(k-1)!-\lambda J_{k-1}(\lambda).
\]
Therefore
\begin{equation}\label{eq:matched-rational-moment}
 \int_0^\infty\frac{e^{-t}V(t)}{t+\lambda}\,dt
 =S_\lambda(V)+V(-\lambda)e^\lambda E_1(\lambda).
\end{equation}
Equations \eqref{eq:matched-condition}--\eqref{eq:matched-rational-moment}
and
\(\Ein(\lambda)=\gamma+\log\lambda+E_1(\lambda)\) give the exact identity
\begin{equation}\label{eq:matched-affine-ein}
 \boxed{
 \mathscr I_i=e^{\lambda_i}R(U_i)+S_{\lambda_i}(V_i)
 -e^{\lambda_i}\beta(U_i)\Ein(\lambda_i).}
\end{equation}
The constant term and the nonzero coefficient of \(\Ein(\lambda_i)\) lie
in \(\Eexp\).  Thus the displayed values are obtained from the distinct
\(\Ein(\lambda_i)\) by an invertible diagonal affine transformation over
\(\Eexp\).  Theorem~\ref{thm:pure-exp-base} proves the claim.
\end{proof}

There is also an exact relation-ideal version.  For a fixed \(\lambda_i\),
let \((U_{ij},V_{ij})\), \(1\leq j\leq r_i\), satisfy
\eqref{eq:matched-condition}, and put
\[
 a_{ij}=e^{\lambda_i}\beta(U_{ij}),
 \qquad
 c_{ij}=e^{\lambda_i}R(U_{ij})+S_{\lambda_i}(V_{ij}).
\]
If \(T_{ij}\) is the coordinate corresponding to
\(\mathscr I_{ij}\), then the complete relation ideal over \(\Eexp\) is
generated by
\begin{equation}\label{eq:matched-full-ideal}
 a_{i1}(T_{ij}-c_{ij})
 -a_{ij}(T_{i1}-c_{i1}),
 \qquad
 1\leq i\leq s,\quad 2\leq j\leq r_i.
\end{equation}
Indeed, \(\mathscr I_{ij}-c_{ij}=-a_{ij}\Ein(\lambda_i)\), and the
\(s\) remaining \(\Ein\)-coordinates are algebraically independent.

The smallest specialization already gives a countable independent family.
Taking \(U=1\) and \(V=-1\) yields
\begin{equation}\label{eq:minimal-matched-integral}
 \boxed{
 \mathscr I(\lambda)
 =e^\lambda\int_0^\infty e^{-t}\log(t/\lambda)\,dt
 -\int_0^\infty\frac{e^{-t}}{t+\lambda}\,dt
 =-e^\lambda\Ein(\lambda).}
\end{equation}
Hence
\(\{\mathscr I(\lambda):\lambda\in\Qbar_{>0}\}\) is algebraically
independent over \(\Eexp\), in the finite-subset sense.

\begin{corollary}[Filtered Touchard integrals]
\label{cor:filtered-touchard-transcendence}
For every \(n\geq0\),
\begin{equation}\label{eq:paired-touchard}
 \mathscr R_n-e\mathsf B_n
 \int_0^\infty e^{-t}\log t\,dt
 =e\mathsf B_n\Ein(1)-\mathsf G_n
\end{equation}
is transcendental over \(\Eexp\).

More generally, let \(P\in\Z[X]\), and put
\[
 q_n=P(E)\mathsf B_n,
 \qquad p_n=P(E)\mathsf G_n.
\]
If \(q_n\ne0\), then
\begin{equation}\label{eq:paired-filtered-touchard}
 \int_0^\infty
 \frac{e^{-x}[P(E)\tau_n](x)}{1+x}\,dx
 -eq_n\int_0^\infty e^{-t}\log t\,dt
 =q_ne\Ein(1)-p_n
\end{equation}
is transcendental over \(\Eexp\).
If \(P\) is monic, then \(q_n>0\) and \(q_n\to\infty\) for all
sufficiently large \(n\).
\end{corollary}

\begin{proof}
The identities follow from \eqref{eq:touchard-remainder} and
\(\int_0^\infty e^{-t}\log t\,dt=-\gamma\).  Their right sides are
nonconstant affine expressions in \(\Ein(1)\) over \(\Eexp\), so they are
transcendental.

For the final assertion, Dobi\'nski's formula gives
\[
 q_n=e^{-1}\sum_{k\geq0}\frac{k^nP(k)}{k!}.
\]
A monic real polynomial is positive at all sufficiently large nonnegative
integers.  Its negative terms in the displayed sum therefore have finite
support.  Choose a positive integer \(L\) beyond that support with
\(P(L)>0\).  The single positive term at \(L\) eventually dominates the
finite negative sum, whose bases are all smaller than \(L\).  It follows
that \(q_n>0\) eventually; choosing \(L\geq2\) also gives \(q_n\to\infty\).
\end{proof}

In particular, for the polynomial \(C_M\) of
Theorem~\ref{thm:bell-gould-arbitrary-modulus}, the two integers \(q_n,p_n\)
are divisible by \(M\).  Dividing
\eqref{eq:paired-filtered-touchard} by \(M\) therefore gives, for all
sufficiently large \(n\), an explicit normalized mixed integral which is
transcendental over \(\Eexp\) and whose \(\Ein(1)\)-coefficient tends to
infinity.  This closes nonvanishing and coefficient growth in the matched
\(e\gamma+\delta=e\Ein(1)\) direction.  It does not prove nonvanishing or
smallness for the unpaired \(\delta\)-linear form and therefore gives no
irrationality or transcendence theorem for \(\delta\) itself.
